% Generated by mceg equations; edit the contract, then regenerate.
% Requires amsmath. Authored math is retained; TeX is not executed here.
% Definitions and hashes do not establish scientific validity.

% equation-sha256: cf21f4682154d64487209033afdbf4fd20369c23cee43fc44f528c6840ebd87f
\paragraph{Equation local\_population\_map (revision 1)}
A Parker velocity and mass continuity supply each layer density; the extended local atom solves charge and statistical equilibrium. Higher shells feed the low states but are not counted in n3. Implicit derivatives include electron-dependent rates. R denotes the complete independent residual, including the conservation closure; its state Jacobian is square.
\begin{equation}
\begin{aligned}
n_{\rm H}(r)=\frac{\dot M}{4\pi r^2m_{\rm H}v(r)},\quad \mathcal R(\boldsymbol n,n_e;T,J)=0,\quad n_e=n_p=n_{\rm H}-\sum_i n_i,\\ n_2=n_{2s}+n_{2p},\quad n_3=n_{3s}+n_{3p}+n_{3d},\quad \frac{d\boldsymbol x}{d\theta}=-\mathcal R_{\boldsymbol x}^{-1}\mathcal R_\theta
\end{aligned}
\label{eq:escape.nlte_5f_spectrum:local_5f_population_5f_map:1}
\end{equation}
\noindent\textbf{Scope:} The wind residual and atomic rates have separate wind and enlarged-atom contracts; x denotes the local solved state, not a new explicit time integrator.

\begin{description}
\item[$n_{\rm H},n_i,n_e,n_p$] Nucleus and particle number densities (cm\textasciicircum{}-3; ID: density).
\item[$r$] Radial coordinate (cm; ID: radius).
\item[$\dot M$] Hydrogen mass loss rate (g s\textasciicircum{}-1; ID: mass\_loss).
\item[$v$] Radial wind speed (cm s\textasciicircum{}-1; ID: speed).
\item[$\boldsymbol x$] Solved local state (model coordinates; ID: state).
\item[$\theta$] Differentiated physical parameters (parameter dependent; ID: parameters).
\end{description}

\noindent\textbf{Assumptions:}
\begin{itemize}
\item Local statistical equilibrium with prescribed radiation and temperature; constant wind molecular weight remains external.
\item The prescribed temperature rises exponentially from its base with scale 0.12 Rp. Irradiation is an unattenuated diluted Planck field. Continuum integration uses Gauss-Legendre nodes in log frequency between ionization thresholds, truncated at 100 eV.
\end{itemize}

% equation-sha256: 65190ee7196dc29d9f7a2271015be3d814ad5d74576b36f1890a05dce5f21e54
\paragraph{Equation branch\_profiles (revision 1)}
Sum branch extinction and emission before forming the source. Each Voigt profile uses full upper/lower radiative lifetimes and thermal Doppler broadening; normalization is over the infinite wavenumber domain, not the truncated plotting grid.
\begin{equation}
\begin{aligned}
L_i=\sum_{j<i}A_{ij}+\delta_{i,2s}A_{2\gamma},\quad \gamma_{ul}=\frac{L_u+L_l}{4\pi c},\quad\int\phi_{ul}(\widetilde\nu)\,d\widetilde\nu=1,\\ \alpha_{\widetilde\nu}=\sum_{ul}\frac{A_{ul}}{8\pi c\widetilde\nu_0^2} \left(\frac{g_u}{g_l}n_l-n_u\right)\phi_{ul},\quad j_{\widetilde\nu}=\sum_{ul}\frac{hc\widetilde\nu_0A_{ul}}{4\pi}n_u\phi_{ul}
\end{aligned}
\label{eq:escape.nlte_5f_spectrum:branch_5f_profiles:1}
\end{equation}
\noindent\textbf{Scope:} H-alpha 3s-to-2p, 3p-to-2s and 3d-to-2p at one common NIST center. Coefficients follow the ExoJAX wavenumber convention.

\begin{description}
\item[$L_i$] Total spontaneous level depopulation rate (s\textasciicircum{}-1; ID: lifetime\_rate).
\item[$\gamma_{ul}$] Natural Lorentz half width (cm\textasciicircum{}-1; ID: gamma).
\item[$\widetilde\nu$] Wavenumber (cm\textasciicircum{}-1; ID: wavenumber).
\item[$\phi_{ul}$] Unit area wavenumber profile (cm; ID: profile).
\item[$\alpha_{\widetilde\nu}$] Net absorption including stimulated emission (cm\textasciicircum{}-1; ID: extinction).
\item[$j_{\widetilde\nu}$] Emissivity per unit wavenumber and solid angle (erg s\textasciicircum{}-1 cm\textasciicircum{}-2 sr\textasciicircum{}-1; ID: emission).
\item[$n_i$] Actual resolved state population (cm\textasciicircum{}-3; ID: bound).
\end{description}

\noindent\textbf{Assumptions:}
\begin{itemize}
\item No pressure broadening or branch-center splitting; no clipping of an inverted population.
\end{itemize}

% equation-sha256: f959f48e0cfa91b6d4ef6e0933e12514170c35312d03c1816fc1cd6c22ce07cd
\paragraph{Equation observer\_transfer (revision 1)}
Each surface patch projects radial, rotation and zonal velocities into its Doppler-shifted profiles. Existing ExoJAX intensity transfer integrates the layered column; visible disk weights and a Gaussian instrument kernel produce line contrast.
\begin{equation}
\begin{aligned}
S_{\widetilde\nu}=j_{\widetilde\nu}/\alpha_{\widetilde\nu},\quad I^{\rm out}_{\widetilde\nu}=I^{\rm bot}_{\widetilde\nu}e^{-\tau_{\rm tot}/\mu} +\sum_k S_{k,\widetilde\nu}e^{-\tau_{<k}/\mu}(1-e^{-\Delta\tau_k/\mu}),\\ \Delta\tau_k=\alpha_{k,\widetilde\nu}\Delta r_k,\quad f_{\widetilde\nu}=\left(\frac{R_p}{R_*}\right)^2 \frac{\mathcal G_R*\sum_q w_q(I^{\rm out}_{q,\widetilde\nu}-I^{\rm bot}_{\widetilde\nu})}{F_{*,\widetilde\nu}}
\end{aligned}
\label{eq:escape.nlte_5f_spectrum:observer_5f_transfer:1}
\end{equation}
\noindent\textbf{Scope:} The NLTE observer requires strictly positive total extinction and rejects zero opacity even if emission is zero. Only the legacy known-source wrapper admits a vacuum and preserves its zero-column derivative. Emitting zero-opacity layers and net gain are unsupported.

\begin{description}
\item[$I$] Specific intensity per unit wavenumber (erg s\textasciicircum{}-1 cm\textasciicircum{}-1 sr\textasciicircum{}-1; ID: intensity).
\item[$S$] Local source function (erg s\textasciicircum{}-1 cm\textasciicircum{}-1 sr\textasciicircum{}-1; ID: source).
\item[$\tau$] Vertical optical depth (1; ID: optical\_depth).
\item[$\mu$] Ray cosine relative to the patch normal (1; ID: cosine).
\item[$f$] Disk integrated planet-to-star line excess (1; ID: contrast).
\end{description}

\noindent\textbf{Assumptions:}
\begin{itemize}
\item Observer transport does not feed radiation back into the local atom.
\end{itemize}

% equation-sha256: d1247fa41e9fb62dde00ebd65002e6e0f5a19e7fcea44faadf2371f76ba0cc04
\paragraph{Equation resolution\_criterion (revision 1)}
Refine one numerical axis at a time and compare at all 87 supplied observation wavelengths. The final candidate also refines every axis together. The same 0.1-uncertainty criterion is fixed before each comparison; failed grids are retained and finer candidates are tested separately.
\begin{equation}
\begin{aligned}
\epsilon_a=\max_{p}\frac{|f_p^{(a\ {\rm refined})}-f_p^{({\rm base})}|}{\sigma_p}<0.1
\end{aligned}
\label{eq:escape.nlte_5f_spectrum:resolution_5f_criterion:1}
\end{equation}
\noindent\textbf{Scope:} N=15 to 20; double layers, disk angles, wavenumber intervals, continuum quadrature and wind collocation intervals independently, then combine these refinements for the final candidate. The initial 32- and 64-layer attempts did not include the joint control. Supplied uncertainties are a numerical scale whose confidence interpretation is unverified.

\begin{description}
\item[$\epsilon_a$] Largest change on refinement of axis a (1; ID: error).
\item[$f_p$] Interpolated spectral contrast at observer pixel p (1; ID: contrast).
\item[$\sigma_p$] Supplied observational uncertainty (1; ID: uncertainty).
\end{description}
